
%\needspace{4\baselineskip}
\color{uniblue}{\subsection{\raggedright{Контроллер присоединения (КП)}}}
\color{black}

\begin{enumerate}

%%%%%%%%%%%%%%%%%%%%%%%%%%%%%%%%%%%%%%%%%%%%%%%%%%%%%%%%%%%%%%%%%%%%%%  ОБЩИЕ СВЕДЕНИЯ %%%%%%%%%%%%%%%%%%%%%%%%%%%%%%%%%%%%%%%%%%%%%%%%%%

\item Общие сведения

\begin{enumerate}

\item 
Контроллер присоединения обеспечивает возможность оперативного управления коммутационными аппаратами с верхнего уровня АСУ ТП, от телемеханики и по месту. Также контроллер присоединения отслеживает состояние управляемых коммутационных аппаратов и реализует оперативную блокировку коммутационных аппаратов в составе присоединения.

\item
Контроллер присоединения предназначен:
\begin{itemize}
\item для оперативного управления коммутационными аппаратами; 
\item для формирования сигнала положения коммутационного аппарата для целей АСУ ТП; 
\item для реализации оперативной блокировки управляемых коммутационных аппаратов.
\end{itemize}

\item
В устройстве контроллер присоединения представлен функциональным блоком «КП», который состоит из одной функции управления выключателем «УВ».

\item
Ввод контроллера присоединения в работу на этапе параметрирования устройства осуществляется программным переключателем SGF1 <<Ввод\_функции>>. Введенное состояние контроллера присоединения соответствует активному состоянию сигнала <<КП: Ввод>>.

\item
Функциональный блок <<КП>> может быть выведен из работы оперативно путем активации сигнала <<ОВ КП>>, о чем сигнализирует активное состояние сигнала <<КП: Оперативный вывод>>.

\item
Логика управления режимами работы функционального блока <<КП>> выполнена в соответствии с рисунком \ref{kp:img_gen}. В таблице \ref{kp:tbl_gen} приведены параметры, необходимые для настройки КП. 

\vspace{2mm}
\begin{figure}[H]
\centering
\includegraphics[width=1\textwidth,height=1\textheight,keepaspectratio]{\fbpath/45. КП В/_latex/img/1.pdf}
\captionof{figure}{Функциональная схема управления режимами работы контроллера присоединения}
\label{kp:img_gen}
\end{figure}

\end{enumerate}
%%%%%%%%%%%%%%%%%%%%%%%%%%%%%%%%%%%%%%%%%%%%%%%%%%%%%%%%%%%%%%%%%%%%%%  УПРАВЛЕНИЕ ВЫКЛЮЧАТЕЛЕМ %%%%%%%%%%%%%%%%%%%%%%%%%%%%%%%%%%%%%%%%%%%%%%%%%%

\item Управление выключателем (УВ)

\begin{enumerate}

\item
В функциональном блоке «КП» управление силовым выключателем представлено функцией «УВ».

\item
Функция формирует команды на отключение и включение выключателя при оперативном управлении.

\item
Местное управление коммутационным аппаратом осуществляется только посредством кнопок \tikz[baseline={(0,-0.00ex)}]{\draw[line width=1.6pt] (0,0) rectangle (1em,1em);\node[anchor=center] at (0.5em,0.5em) {\textbf{I}};}~–~<<Включить>> и \tikz[baseline={(0,-0.00ex)}]{\draw[line width=1.6pt] (0,0) rectangle (1em,1em);\node[anchor=center] at (0.5em,0.5em) {\textbf{O}};}~–~<<Отключить>>. При этом необходимо перевести устройство в режим <<Местное>> с помощью заданной функциональной кнопки. В данном случае управление по локальной сети и с помощью выносных кнопок или ключей будет невозможно.
\item
Дистанционное управление может реализоваться через дискретные входы, которые сконфигурированы как <<Включить В от ТУ>> и <<Отключить В от ТУ>>, а также по локальной сети. В дистанционном режиме управления команды по локальной сети предусмотрены постоянно.
\item
Управление выключателем также осуществляется с помощью пульта управления (ПУ), который предназначен для формирования команд управления выключателем путем автоматической подачи на катушки электромагнитов электрического напряжения. Управление с ПУ может реализоваться через дискретные входы, которые сконфигурированы как <<Включить В от ПУ>> и <<Отключить В от ПУ>>.

\item
Команда оперативного отключения <<УВ: Отключить>> формируется в следующих случаях:
\begin{itemize}
	\item при появлении команды на оперативное отключение выключателя (<<Отключить В~от~ПУ>> или <<Откл. В~ИЧМ>>) и выбранным управлением по месту (сигнал <<Местное>> находится в активном состоянии);
	\item при появлении команды на оперативное отключение выключателя (<<Отключить В~от~ТУ>> или <<Откл. В~АСУ>>) и выбранным дистанционным управлением (сигнал <<Дистанционное>> находится в активном состоянии).
\end{itemize}
\item
Прохождение команды на отключение блокируется при наличии следующих сигналов:
\begin{itemize}
	\item запрет отключения от функции контроля силового выключателя <<КСВ: Блокировка откл>>;
	\item отключенное положение силового выключателя <<В отключен (б/к)>> (с контролем сигнала <<В~включен~(б/к)>>);
	\item превышение времени переключения силового выключателя <<УВ: Прев. времени~пер.>>;
	\item местное управление силовым выключателем (положение переключателя <<М/Д>> в приводе выключателя) <<Ключ М/Д привода В>>.
\end{itemize}

\item
Команда оперативного включения <<УВ: Включить>> формируется в следующих случаях:
\begin{itemize}
	\item при появлении команды на оперативное включение выключателя (<<Включить В~от~ПУ>> или <<Вкл. В~ИЧМ>>) и выбранным управлением по месту (сигнал <<Местное>> находится в активном состоянии);
	\item при появлении команды на оперативное включение выключателя (<<Включить В~от~ТУ>> или <<Вкл. В~АСУ>>) и выбранным дистанционным управлением (сигнал <<Дистанционное>> находится в активном состоянии).

\end{itemize}
\item
Прохождение команды на включение блокируется при наличии следующих сигналов:
\begin{itemize}
	\item запрет включения от функции контроля силового выключателя <<КСВ: Блокировка вкл>>;
	\item включенное положение силового выключателя <<В~включен (б/к)>> (с контролем сигнала <<В~отключен~(б/к)>>);
	\item превышение времени переключения силового выключателя <<УВ: Прев. времени~пер.>>;
	\item местное управление силовым выключателем (положение переключателя <<М/Д>> в приводе выключателя) <<Ключ М/Д привода В>>;
	\item появление сигнала аварийного отключения силового выключателя <<КСВ: В аварийно отключен>> (прохождение сигнала определяется состоянием программного переключателя SGF2 <<Бл\_вкл\_от\_авар\_откл>>).
\end{itemize}

\item
При наличии сигналов <<УВ: Включить>> и <<УВ: Отключить>> длительностью больше, чем время уставки <<T1>> формируется сигнал превышения допустимого времени переключения -- <<УВ: Прев. времени пер.>>. Активное состояние операции включения или отключения сигнализируется выходным сигналом <<УВ: Идет переключение>>.

\item
Логика функции формирует сигналы состояния выключателя <<УВ: Отключено>>, <<УВ: Включено>>, <<УВ: Не определено>> и <<УВ: Неисправность>>, при этом сигналы <<УВ: Неисправность>> и <<УВ: Не определено>> выдаются с выдержкой времени, задаваемой уставкой Т2 <<Тблк>>.

\item
Ввод функции в работу на этапе параметрирования устройства осуществляется программным переключателем SGF1 <<Ввод\_функции>>. Информация о введенной функции отображается выходным сигналом <<УВ: Ввод>>.

\item
Алгоритм работы функции выполнен в соответствии с рисунком \ref{kp:img_uv}.

\vspace{3mm}
\begin{figure}[H]
\centering
\includegraphics[width=0.9\textwidth,height=0.9\textheight,keepaspectratio]{\fbpath/45. КП В/_latex/img/uv.pdf}
\captionof{figure}{Функциональная схема логики управления выключателем}
\label{kp:img_uv}
\end{figure}

\end{enumerate}

\end{enumerate}
\FloatBarrier % Все float'ы ДО этой строки будут размещены здесь.